\subsection{Generic infinite-support nondegeneracy}
Let $X_A$ be the real Banach space obtained by adjoining the constant Fourier
mode to the two weighted moments in \eqref{eq:wiener-moments}:
\[
 \|f\|_{X_A}=|\widehat f(0)|+M_0(f)+M_A(f).
\]
On absolutely summable Fourier sequences define the orbit-sum map
\[
 (\mathcal Ph)_{\mathcal O}=\sum_{k\in\mathcal O}\widehat h(k).
\]
Because the oriented frequency orbits partition $\mathbb Z^2\setminus\{0\}$,
\begin{equation}\label{eq:wiener-orbit-cont}
 \|\mathcal Ph\|_{\ell^2(\mathcal O)}
 \leq\|\mathcal Ph\|_{\ell^1(\mathcal O)}
 \leq\|\widehat h\|_{\ell^1}.
\end{equation}
Moreover
\[
 h_g-h_f=B_A(g-f,g)+B_A(f,g-f),
\]
so \eqref{eq:wiener-bilinear} yields
\begin{equation}\label{eq:wiener-gamma-cont}
 \|\widehat{h_g-h_f}\|_{\ell^1}
 \leq\pi\{M_0(g-f)M_A(g)+M_0(f)M_A(g-f)\}.
\end{equation}
Thus $g\mapsto\mathcal Ph_g$ is continuous from $X_A$ to
$\ell^2(\mathcal O)$ and
$\Gamma_A(g)=\|\mathcal Ph_g\|_2^2$ is continuous.

\begin{theorem}[Generic weighted-Wiener nondegeneracy]
\label{thm:wiener-generic}
Modulo constants, the set
\[
 \{g\in X_A:\Gamma_A(g)>0\}
\]
is open and dense.  Its exact-cancellation complement is closed and nowhere
dense.
\end{theorem}
\begin{proof}
Openness follows from continuity.  For density, let $p_R$ be the Fourier
truncation of $g$ to $\|v\|_1\leq R$.  Absolute convergence of the weighted
moments gives $p_R\to g$ in $X_A$.  If $g$ is nonconstant then $p_R$ is
nonconstant for all sufficiently large $R$, and the unrestricted
finite-Fourier theorem gives $\Gamma_A(p_R)>0$.  Constants are approximated by
adding an arbitrarily small nonzero cosine mode.  Hence the positive set is
dense; its closed complement has empty interior.
\end{proof}

For $\rho>0$ let
\[
 X_\rho=\left\{f:\sum_v e^{\rho\|v\|_1}|\widehat f(v)|<\infty\right\}.
\]
Since
\[
 M_0(f)\leq C_\rho\|f\|_{X_\rho},\qquad
 M_A(f)\leq\|A\|_{1\to1}C_\rho\|f\|_{X_\rho},\qquad
 C_\rho=\sup_{n\ge1}ne^{-\rho n}<\infty,
\]
the embedding $X_\rho\hookrightarrow X_A$ is continuous and finite
trigonometric polynomials are dense in $X_\rho$.  Therefore
$\{\Gamma_A>0\}$ is open and dense in every fixed analytic Wiener strip as
well.

There is also a natural Sobolev version.  For $s>2$, Cauchy--Schwarz gives
\[
 M_0(f)\leq\|f\|_{H^s}
 \left(\sum_{v\in\mathbb Z^2}\|v\|_1^2\langle v\rangle^{-2s}\right)^{1/2},
 \qquad M_A(f)\leq\|A\|_{1\to1}M_0(f),
\]
and the lattice sum converges because its radial shell is
$O(n^{3-2s})$.  Hence $H^s(\mathbb T^2)$ embeds continuously in $X_A$ and
trigonometric polynomials are dense there.

\begin{corollary}[Sobolev and smooth genericity]
For every $s>2$, modulo constants,
\[
 \{g\in H^s(\mathbb T^2;\mathbb R):\Gamma_A(g)>0\}
\]
is open and dense; the zero-variance set is closed and nowhere dense.  The
same conclusion holds in $C^\infty(\mathbb T^2)$ with its Fr\'echet topology.
\end{corollary}

Thus the known smooth exact-cancellation examples are topologically
exceptional even though they remain valid.  The separated-support example is
compactly supported Gevrey-$s$ with $s>1$, not analytic, while the nontrivial
entire temporal all-pass gauges of one harmonic are proved elsewhere to lie
in no finite Gevrey class.  An all-members theorem for arbitrary nonconstant
analytic observables, or an analytic exact-cancellation counterexample,
remains open.
